% Fresh proof, 2026-10-22. Frozen reconstructed inputs are unchanged.
\section{R2 with bounded \(a\), or an infinitesimal paired-event parameter}
\label{sec:r2-shifted-local}

\paragraph{Exact statement.}
Use the actual original weak-C2 source and all its original strict
gaps, with \(d=1\), \(\lambda=\beta=0\), and \(n\geq0\).
Suppose \(t_s\) is positive unbounded and
\(B_s,a_sB_s,t_sB_s\) are positive infinitesimals.
Put \(e=1-r\), \(R_s=1+ha_sB_s\).  On the source subset
\[
       a_s=O(1)\quad\hbox{or}\quad
       p=a_s e/R_s\in\mathrm I_c,                     \tag{RL1}
\]
the entire parent is separated tensor SOS for every independent
projective testing pair and every real or complex state.
Literal \(n=0\) and every opening of \(r,h,1-h,g\) are retained.
For the all-testing conclusion, the local proof below is joined
explicitly to R2O1--R2O24.

The two fields and the bounded-Gram convention are those of EN1--EN4.
In particular no mixed inverse is used.  Every coefficient bound
below refers to a finite array in \(B_c\otimes B_t\); the rows need
not be bounded.

\subsection{Actual source entry and bounds}

Use the exact compact source dictionary RSR2--RSR8, with
\[
 j=a_sc_s/R_s,\quad B=B_sR_s/c_s,\quad T=t_sc_s/R_s,\quad
 p=a_se/R_s,\quad H=(1-h)/(t_sc_s),\quad c_s=1+a_sB_sr.
\]
Set
\[
 S=1/T,\quad \ell=HT,\quad \bar h=1-\ell,\quad
 x=jB,\quad y=TB,\quad B=Sy,\quad q_z=j/S .
                                                               \tag{RL2}
\]
Thus \(S,x,y\) are positive infinitesimals,
\(0<p<\min(1,j)\), \(0<\ell<1\), and \(j\asymp a_s\).
Both \(\ell\) and \(\bar h\) may be arbitrarily small.
The original physicality inequality RER2--RER7 gives \(jy\in\mathrm I_c\).
In terms of the actual finite polynomials \(\mathcal P,\mathcal L,
\mathcal N,\mathcal F,\mathcal K\) of RSR, put
\[
 K=j+1-p,\quad k=K/(j+1),\quad
 F_b=j+1+p(1-p)/j,\quad E=1+B(\ell+p-1).
                                                               \tag{RL3}
\]
Here \(1/2<k<1\) and \(1\le F_b/(j+1)<2\).  Direct substitution
\(B=Sy,T=1/S,H=S\ell\) in those finite polynomials gives
\[
 \mathcal P/(j+1)\sim1,\qquad
 S\mathcal L/K\sim1,\qquad S\mathcal F/F_b\sim1.         \tag{RL4}
\]
For clarity, the terms after division contain only bounded
\(p/j,j/(j+1),K/(j+1)\), multiplied by \(S,y\), or \(jy\).
No inverse of \(\ell\), \(\bar h\), or \(j-p\) is involved.

The same substitution in the exact numerator of \(n\) yields
\[
 \mathcal N=p(1+\epsilon_p)
                -yj(j+1)(1+\epsilon_y),\qquad
 \epsilon_p,\epsilon_y\in\mathrm I_c.                  \tag{RL5}
\]
For an explicit finite division certificate, the zero-\(p\) term is
\[
 \mathcal N|_{p=0}
 =-yj(j+1)\{1+S(1-\ell)\}
             \{1+S\ell yj/(j+1)\}.
\]
The quotient \((\mathcal N-\mathcal N|_{p=0})/p\) is a
finite polynomial.  This grouping can be read termwise from RSR5:
the first
coefficient is \(1+y(j+1)+O(Sy(j+1)+Sy^2(j+1))\);
the second coefficient is \(1+O(S+y+Syj)\).
Using \(\mathcal N\geq0\), therefore,
\[
 \eta_2=\frac{yj(j+1)}p=O(1),\qquad
 \eta_2\in\mathrm I_c\ \Longrightarrow\
 n\sim\frac{p}{2(j-p)(j+1)}>0 .                        \tag{RL6}
\]
The estimate includes the face \(\mathcal N=0\), which can occur
only in the positive-unit \(\eta_2\) case.

Let \(\mu=m(1-p)\) and \(h_0=\mu+1/2\).
The complete source formulas now imply
\[
\begin{gathered}
 \mu\sim\frac{jp}{2(j+1)},\quad h_0\asymp1,\quad
 CH\sim k/2,\quad SA\bar h\sim k/2,\quad D/S\sim k/2,\\
 \alpha\sim\frac{Sy^2K^2}{2(j+1)},\quad
 f_c\sim\frac{j^2y^2F_b}{2(j+1)},\quad
 r_c\sim\frac{jy^2(1-p)K}{2(j+1)} .
\end{gathered}                                                   \tag{RL7}
\]
These use the exact masses in RFN12 and the closed RSR formulas.
All displayed positive masses are strictly positive.
Write \(\kappa=-k_c>0\).  The exact positive-surplus formula
RSR19 proves this sign.  Its substituted finite polynomial gives
\[
 \frac{\kappa}{r_c}=O(p+Sy)\quad(p\in\mathrm I_c),\qquad
 \frac{\kappa}{jy^2}=O(1).                              \tag{RL8}
\]
For the remaining case in (RL1), \(p\) and \(j\) are positive
bounded units, and
\[
 \frac{\kappa}{f_c}\sim\frac{pK}{jF_b}\leq\frac12,
 \qquad jF_b-2pK=(j-p)K\geq0.                          \tag{RL9}
\]
The inequality is a bound on real residues when combined with
the preceding asymptotic equality.  On this case \(\eta_2\) is
infinitesimal, and the effective current defined below is a
positive bounded unit.

\subsection{The entire parent and its harmonic shift}

Retain the complete original quarter-parent RFN4.  Its exact inverse
state change at \(w=(W+uY)/(1+u^2)=1\) is
\[
 X=N_0+pv-uz,\quad Z=z+uN_0+puv,\quad
 W=1+u(M_0+Hvz),\quad Y=u-M_0-Hvz.                    \tag{RL10}
\]
Here \(z>0\) is arbitrary.  Retain \(mpD_u^2N_0^2\) and divide
the remaining whole parent by \(D_u=1+u^2\).
The exact expression is RFN7, with both complete state normals,
hard entries RFN9, and scalar RFN12.

Use the additional exact change
\[
 q=vz-Tu,\qquad L=M_0+Hq,\qquad
 M_0=L-Hq,\qquad L=\sqrt{S\ell\bar h}\,\widehat L .
                                                               \tag{RL11}
\]
All divisions are pure source or pure testing.  In particular
this change does not invert the harmonic mixed unit
\(1+\rho u^2\).
The complete diagonal of \(\widehat L\) is
\[
 S\ell\bar h\{D_u(CD_v+D+\tfrac12)+AD_v+
                   (\alpha+\tfrac12)v^2\},
                                                               \tag{RL12}
\]
with positive real residue \(\overline k/2\) and bounded
coefficient array through both openings.

Use the circuit coordinates \(U,V,Z_*,Q_*,B_*,R_*,\Psi\)
of SH1--SH2, and \(b=v-uz/j=\sqrt y B_*\).
Before spending any scalar \(n\)-mass, complete the exact correlated
block, as in RFN13--RFN16:
\[
 h_0N_0^2+n(N_0+jb)^2
 =(n+h_0)(N_0+c_bb)^2+M_e b^2,\quad
 c_b=\frac{nj}{n+h_0},\quad M_e=\frac{nh_0j^2}{n+h_0}.  \tag{RL13}
\]
Keep \(nv^2(N_0+pb)^2\) and all three additional positive hard
squares RFN15 whole, including after the change
\(\widehat N=\sqrt{n+h_0}(N_0+c_bb)\).
At \(n=0\), put \(c_b=M_e=0\) and omit the current row.
For \(n>0\), \(c_b^2/M_e=n/[h_0(n+h_0)]\) is bounded.

\paragraph{The exact shifted scalar.}
Let \(F\) be RFN12.  The shifted scalar is
\(F-2HqL_M+H^2q^2H_M\).
Use the whole \(d_cV^2Q_*^2\) term in its \(q^2\) diagonal.
Using the exact actual-source identity \(S_0=0\) in RFN22,
the resulting coefficient on \((\sqrt S q)^2\) is exactly
\[
 D_q=S\{A\bar h^2+\ell^2(C+D+\tfrac12)D_u\}
       -\frac{f_c}{Tx}-TH^2C\,u^2v^2
       -\frac{k_cv^2}{yTx}.                            \tag{RL14}
\]
It has bounded array and positive real residue \(\overline k/2\).
Indeed \(f_c/(Tx)=O(jy)\) is infinitesimal, \(TH^2C=\ell CH\)
is bounded, and \(\kappa/(yTx)=\kappa/(jy^2)\) is bounded.
Its equivalent \(Q_*^2\) mass has order \(jy\).
The only extra scalar cross terms are
\[
 -2Hq\{\mu uvb+v(\alpha-yCHu^2)R\},\qquad
 R=z-uv/(Sy).                                        \tag{RL15}
\]
The budgets below allocate their entire expressions.

The remaining signed scalar terms are exactly
\(e_0V^2B_*^2+k_cZ_*^2(V^2-1)^2\), where
\(e_0=p^2y^2/2-g_c\).
The negative part of \(e_0\) is at most \(g_c\).
SH2 purchases it at relative circuit cost
\[
 \frac{g_c}{f_c}(1+V^2)
   =O\bigl(S(y+v^2)\bigr)\in\mathrm I.                \tag{RL16}
\]
Use \(Z_*(V^2-1)=-VQ_*-R_*\) for the other signed term.
If \(p\) is infinitesimal, (RL8) pays its opposite part
directly from \(r_cR_*^2\), and its \(VQ_*\) cost is
\(O(v^2)\).
Otherwise \(p,j,M_e\) are positive units.  Young's exact
square identity with parameter \(1/6\), followed by SH1 at
\(a=3/2\), gives the legal SOS domination
\[
 \kappa Z_*^2(V^2-1)^2
 \preceq_{\rm sos}
 \frac74\kappa\Psi+
 \frac{21}{2}\kappa V^2Q_*^2+
 \frac72\kappa V^2B_*^2 .                             \tag{RL17}
\]
The circuit cost is at most \(7/8+o(1)\), by (RL9).
The two other relative costs are \(O(v^2)\).
Thus fixed positive real fractions of \(D_q(\sqrt S q)^2\),
\(r_cR_*^2\), \(f_c\Psi\), and \(M_eb^2\) remain;
the last is absent only at \(n=0\).
The opposite mass \(r_c\) was not spent in the second case.

\subsection{All helper rows and both complete normals}

As usual, a free row permits any fixed real square pivot at
infinitesimal cost from the preceding retained scalar forms.
SH2 proves that both \(E_u=u(uvz-jy)\) and \(E_v=v(uvz-jy)\)
are free.  The row \(R_u=v(u^2-x)\) is free as well.
For \(p\) infinitesimal, use the actual RFN19 rational Gram:
its three costs are \(O(S^2y)\), \(O(Syu^2)\),
and \(O(Su^2/(1-p))\).
When \(p,j,M_e\) are units, apply SH1 to the \(R_*^2\)
term in RFN19.  The resulting costs are
\[
 O(S^2y+Su^2/j),\quad O(Syu^2+Su^2v^2),\quad
 O(jSu^2v^2/M_e),                                    \tag{RL18}
\]
all infinitesimal.

The rows \(pv^2b\) and \(\sqrt S\,\mu uvb\) are free.
If \(\eta_2\) is infinitesimal, (RL6) and (RL13) imply
\[
 \frac{p^2}{M_e}
  =\frac{p^2}{h_0j^2}+\frac{p^2}{nj^2}=O(1).
                                                               \tag{RL19}
\]
They then cost \(O(v^4)\) and \(O(Su^2v^2)\).
If \(\eta_2\) is a positive unit, the actual source gives
\(p^2y/f_c=O(1)\) and \(\mu=O(p)\).
The second identity in SH2 buys the same rows with costs
\[
 O\bigl(v^2(y+v^2)\bigr),\qquad
 O\bigl(Su^2(y+v^2)\bigr),                             \tag{RL20}
\]
including \(n=0\).

The full opposite term \(\sqrt S\,v(\alpha-yCHu^2)R\)
is free by SH4--SH8.  Also \(j\alpha b\) is free:
\[
 j\alpha b=-(\alpha/B)R_u-\alpha uR,\qquad
 \alpha/B=O(y(j+1))\in\mathrm I_c.
                                                               \tag{RL21}
\]
For \(p\) infinitesimal, the direct \(r_c\) cost for
\(\alpha uR\) is \(O(Sy^2(j+1)^2u^2)\).
In the unit \(p,j,M_e\) case use SH1; the three costs
are \(O(Sy^2u^2)\), \(O(Sy^2u^2v^2)\), and
\(O(Sjy^2u^2v^2/M_e)\).  All are infinitesimal.
These statements account for the whole source \(\alpha\), not
only its leading coefficient.

\paragraph{The first normal.}
The exact RFN21 source cancellations give, after (RL11),
\[
 L_N'=[j(n-\alpha)+p(n+\tfrac12)v^2]b
             +K_sR_u+\kappa_NE_v,                     \tag{RL22}
\]
\[
 K_s=-\mu+(\alpha+\tfrac12)(\ell+p-1),\qquad
 \kappa_N=-(n+\tfrac12)(1-p/j).
\]
Both coefficients are bounded.  For \(\kappa_N\), use the
exact numerator of \(n\) and
\(n(j-p)/j=O(p/[j(j+1)])\); this cancels the \(r\) opening.
After the entire correlated completion (RL13), the remaining
normal is
\(-j\alpha b+\tfrac12pv^2b+K_sR_u+\kappa_NE_v\).
Every row is free by (RL18)--(RL21).
Division by \(\sqrt{n+h_0}\) is bounded.
The shift generates couplings to \(\sqrt{M_e}b\) with bounded
coefficients \(c_b/\sqrt{M_e}\); they are absent at \(n=0\).
Large fixed real helper pivots absorb these couplings by exact
Young identities.

\paragraph{The second normal and the entire hard block.}
The complete RFN22 normal after (RL11) is
\[
 L_L=(C_q-HH_M)q+\mu uvb+v(\alpha-yCHu^2)R,\qquad
 C_q-HH_M=-D-H(C+D+\tfrac12)u^2 .                      \tag{RL23}
\]
The coefficient of \(\sqrt S q\) after the normalization
of \(L\) is infinitesimal; \(D\) is infinitesimal and
\(H(C+D+1/2)\) is bounded.
The other two terms are free by (RL20) and SH6--SH8.
Exactly the same rows pay (RL15), whose first factor is
\(\sqrt S q\) with a bounded source multiplier.

The remaining hard off-diagonal is
\((\alpha+1/2)uv\sqrt{S\ell\bar h}/\sqrt{n+h_0}\),
which is infinitesimal.  Its \(c_b b\) correction has the
same bound after division by \(\sqrt{M_e}\).
The complete finite Gram therefore has a positive definite
real residue, with principal normal pivots
\[
                   1,\quad\overline k/2,\quad\overline k/2
                                                               \tag{RL24}
\]
on \(\widehat N,\widehat L,\sqrt S q\), together with the
retained strict scalar and purchased helper pivots.
All extra positive hard squares and both original retained
current packets are restored whole.  The strict bounded-Gram
lemma proves tensor SOS for the entire parent.

\subsection{State and outer-chart exhaustion}

The fresh full-state sign lemma SJ1--SJ10 permits positive
\(u,v\), and on the remaining negative-core chart gives
\(wz>0\).  Divide by the nonzero pure-testing \(w\) to obtain
the normalization used above with arbitrary \(z>0\).
No estimate in this proof uses its size.  Literal axes and
homogeneous endpoints have the direct original-core SOS exit.
Real and imaginary state parts separate.
The exact outer theorem R2O1--R2O24 covers every other testing
pair on all original R2 sources.  This proves the stated
all-tests conclusion on (RL1).

The complementary source \(a_s\) unbounded and \(p\) a positive
unit is not inferred from this proof.  It requires the separate
TU source comparison and additional-current conductor.
